\subsection{Generalization Across Graph Distributions}
\label{sec:relwork:generalization}
RQ5 trains on one structural distribution and evaluates on another, so its plausibility
rests on what is known about GNNs under such shift.

\subsubsection{Size and Distribution Shift in GNNs}
\label{sec:relwork:generalization:shift}
Message-passing weights are shared across nodes and layers, so a trained model accepts a
graph of any size, a point the inductive line established in principle
\citep{hamilton2017graphsage}. Accepting an input is not generalising to it:
\citet{buffelli2022sizeshift} show accuracy degrading when test graphs are much larger
than training graphs, and propose regularising toward size-invariant representations.
Train-reduced, infer-full is therefore plausible rather than guaranteed, which is
precisely the posture RQ5 takes.

\subsubsection{Train-Test Structural Mismatch}
\label{sec:relwork:generalization:mismatch}
The direct evidence is one positive and one negative result. SCAL \citep{huang2021scal}
trains on coarsened graphs and evaluates on the originals with usable accuracy on
node-level benchmarks. CTS-Bench \citep{cts_bench2026} does the analogous cross-state
evaluation on netlists and lands below $R^2 = 0$
(\ref{sec:relwork:domain:circuits}). Together they make RQ5 an open question rather than
a formality, and they are why the study carries a positive control, the provably lossless
compression of~\ref{sec:method:architecture:exact}, without which a poor cross-state
result could not be told apart from a broken evaluation path (\ref{sec:intro:rq5}).
